\cvheader{\CandidateCityDE}{Agraringenieur | Milchwissenschaftler | Angewandte Nutztierforschung}

\section*{Profil}
% CUSTOMIZE: Die wichtigsten Anforderungen aus der Stellenausschreibung spiegeln.
Agraringenieur und Milchwissenschaftler mit über vier Jahren Erfahrung in angewandter Nutztierforschung, Versuchsplanung, Projektkoordination und statistischer Auswertung. Entwicklung skalierbarer, kostengünstiger Gassensorsysteme für das Emissionsmonitoring in der Nutztierhaltung im Rahmen der Promotion. Erfahren in Studien auf Praxisbetrieben, der Verknüpfung von Herdenmanagement- und Umweltdaten sowie der Aufbereitung wissenschaftlicher Ergebnisse.

\section*{Berufserfahrung}
\entry{Wissenschaftlicher Mitarbeiter (Doktorand)}{März 2023--heute}{Leibniz-Institut für Agrartechnik und Bioökonomie (ATB), Potsdam}{Freie Universität Berlin}
\begin{cvitems}
  \item Planung und Durchführung interdisziplinärer Feldforschung in Milchviehställen sowie Koordination mit Wissenschaftlern, technischen Mitarbeitenden und Praktikanten.
  \item Entwicklung und Validierung skalierbarer Low-Cost-Gassensorsysteme auf Basis von NDIR- und elektrochemischen Sensoren.
  \item Auswertung von Langzeitdaten mit R und Verknüpfung von Emissions- und Stallklimadaten mit Herdeninformationen.
  \item Erstellung von Publikationen, technischer Dokumentation und Konferenzbeiträgen; Betreuung von Praktikanten und Koordination von Feldkampagnen.
\end{cvitems}

\entry{Wissenschaftlicher Mitarbeiter (Vollzeit)}{Feb. 2022--Feb. 2023}{Leibniz-Institut für Agrartechnik und Bioökonomie (ATB), Potsdam}{}
\begin{cvitems}
  \item Planung und Durchführung angewandter Forschung zu Stallklima und gasförmigen Emissionen unter Praxisbedingungen.
  \item Aufbau, Betrieb und Validierung von FTIR-, CRDS-, GC-MS-, TDLAS-, PAS- und Umweltmesssystemen; Auswertung und Dokumentation komplexer Versuchsdaten.
\end{cvitems}

\entry{Wissenschaftliche Hilfskraft / Masterarbeit}{Apr. 2021--Jan. 2022}{Leibniz-Institut für Agrartechnik und Bioökonomie (ATB), Potsdam}{}
\begin{cvitems}
  \item Untersuchung der räumlichen Verteilung von Methan, Ammoniak und Kohlendioxid in einem natürlich belüfteten Milchviehstall.
\end{cvitems}

\section*{Ausbildung}
\entry{Promotion in Agrartechnik (laufend)}{}{Freie Universität Berlin}{}
\entry{M.Sc. Dairy Science, Note 1,6}{}{Christian-Albrechts-Universität zu Kiel}{}
Schwerpunkte: Wiederkäuerernährung, Herdengesundheit, Precision Livestock Farming, Futterqualität, Ökoeffizienz, Milchwirtschaft und biometrische Versuchsplanung.

\entry{B.Tech. Dairy Technology}{2015--2019}{College of Dairy Science and Food Technology, Indien}{}

\section*{Fachprofil \& Nachweise}
\textbf{Wissenschaftliche Leistungen:} 1 peer-reviewte Erstautorenpublikation; 1 Manuskript im Begutachtungsverfahren; 8 Vorträge bei 9 Konferenzen; Mitarbeit an 4 Forschungsprojekten.\par
\textbf{Methoden:} Studienplanung; statistische Analyse (R); wissenschaftliches Schreiben; Literaturrecherche; Projektkoordination; R, Python, MATLAB und LaTeX.\par
\textbf{Milchvieh und F\&E:} Wiederkäuerernährung; Herdendaten; Haltungssysteme; Umweltmonitoring; Sensorentwicklung und Feldvalidierung.

\section*{Engagement, Sprachen \& Praktische Angaben}
\textbf{Engagement:} Co-Chair, EurAgEng Young Professionals Network; Promovierendenvertretung, ATB; Communication Manager, Leibniz PhD Network.\par
\textbf{Sprachen:} Hindi (Muttersprache); Deutsch (B2); Englisch (B2). \quad
\textbf{Arbeitserlaubnis:} Niederlassungserlaubnis; kein Sponsoring erforderlich.\par
\textbf{Führerschein:} Indische Kategorie B und International Driving Permit; deutsche Klasse B in Vorbereitung.\par
\textbf{Referenzen:} Prof. Dr. Thomas Amon und Dr.-Ing. David Janke, ATB.

